
\subsubsection*{2.2.1}
\phantomsection
\label{sga3:VIIB:2.2.1}

\(\mathbf O_k\) 上的双代数 \(\mathbf H\) 称为\emph{平坦的}，若其底模平坦 \textup{（参见 \SGARefLink{sga3:VIIB:1.2.1}{1.2.1}）}。\footnote{\label{sga3:VIIB:nde:79}编者扩充了下面的原文。}若 \(\mathbf H\) 平坦，则由 \SGARefLink{sga3:VIIB:1.3.5}{1.3.5}，\[
  A=\Gamma^*(\mathbf H)
\] 为拓扑平坦的预有限 \(k\)-代数，且 \(\operatorname{Spf}^*(\mathbf H)\) 作为从 \[
  (\operatorname{Alf}/k)^{\mathrm{op}}
  =
  \operatorname{Vaf}/k
\] 到集合范畴的函子，同构于函子 \[
  \operatorname{Spf}(A):
  C\longmapsto
  \operatorname{Hom}_{\operatorname{Vaf}/k}
    \bigl(\operatorname{Spf}(C),\operatorname{Spf}(A)\bigr).
\]。所以 \(\operatorname{Spf}^*(\mathbf H)\) 上的群结构赋予 \[
  \mathscr G(\mathbf H)=\operatorname{Spf}(A)
\] 形式群结构，其具体描述如下。

对 \(\operatorname{Alf}/k\) 的每个对象 \(C\)，\(\mathbf H(C)\) 的底 \(C\)-模为投射模。因此由 \SGARefLink{sga3:VIIB:1.2.3.A}{引理 1.2.3.A}，对 \(n\) 归纳得到自然同构 \[
\begin{aligned}
  \bigl(\mathbf H(C)^*\bigr)^{\widehat\otimes(n+1)}
  &\simeq
  \operatorname{Hom}_C
    \bigl(
      \mathbf H(C),
      \bigl(\mathbf H(C)^*\bigr)^{\widehat\otimes n}
    \bigr)
  \\
  &\simeq
  \operatorname{Hom}_C
    \bigl(
      \mathbf H(C),
      \bigl(\mathbf H(C)^{\otimes n}\bigr)^*
    \bigr)
  \\
  &\simeq
  \bigl(\mathbf H(C)^{\otimes(n+1)}\bigr)^*.
\end{aligned}
\]。对 \(n=1,2\)，由此可知 \(\mathbf H(C)\) 上的 \(C\)-代数结构赋予 \(\mathbf H(C)^*\) 满足 \SGARefLink{sga3:VIIB:2.1}{2.1}\textup{\SGARefLink{sga3:vii:paragraph:5-1:equation:i}{(i)}}--\textup{\SGARefLink{sga3:VIIB:2.1:item:iii}{(iii)}} 的对角映射，且关于 \(C\) 具有函子性。

因此 \[
  A
  =
  \Gamma^*(\mathbf H)
  =
  \varprojlim_{\mathfrak l}\mathbf H(k/\mathfrak l)^*
\] 具有 \(\operatorname{Alp}/k\) 中的余群结构，从而定义所述 \(\mathscr G(\mathbf H)\) 上的形式群结构。

反之，设 \(G\) 为仿射代数是 \(A\) 的拓扑平坦形式 \(k\)-群，以 \[
  \mathbf H(G)=\mathbf V_k^f(A)
\] 表示其 \(\mathbf O_k\)-余代数 \textup{（参见 \SGARefLink{sga3:VIIB:1.2.3}{1.2.3}）}。于是对每个有限长度 \(k\)-代数 \(C\)，对角态射 \[
  \Delta_A:A\longrightarrow A\widehat\otimes_kA
\] 诱导 \(C\)-线性映射 \[
  \mathbf V_k^f(A)(C)
  \otimes_C
  \mathbf V_k^f(A)(C)
  \longrightarrow
  \mathbf V_k^f(A)(C),
\]，使余代数 \(\mathbf V_k^f(A)(C)\) 成为 \(C\)-双代数。称 \(\mathbf H(G)\) 为形式群 \(G\) 的\emph{协变双代数}。\footnote{\label{sga3:VIIB:nde:80}编者添加形容词“协变”，以明确函子 \(G\mapsto\mathbf H(G)\) 为协变函子；[Di73]，I \S~2.14 使用了此术语。}于是由 \SGARefLink{sga3:VIIB:1.3.5.D}{命题 1.3.5.D} 得到下述命题。

\medskip\noindent\textbf{命题。}\phantomsection \label{sga3:VIIB:2.2.1.prop}{\itshape\begin{enumerate}\item[\textup{(i)}]函子 \[
  G\longmapsto\mathbf H(G),
  \qquad
  \mathbf H\longmapsto\mathscr G(\mathbf H)
\] 定义拓扑平坦形式 \(k\)-群范畴与平坦 \(\mathbf O_k\)-双代数范畴之间的等价。\footnote{\label{sga3:VIIB:nde:81}回顾本讲中“双代数”指“余交换 Hopf 代数”。}

\item[\textup{(ii)}]此等价“与基变换交换”：若 \(k\to K\) 为伪紧环态射，则 \[
  \mathbf H\bigl(G\widehat\otimes_kK\bigr)
  =
  \mathbf H(G)\otimes_kK
\]，且 \[
  \mathscr G\bigl(\mathbf H\otimes_kK\bigr)
  =
  \mathscr G(\mathbf H)\widehat\otimes_kK.
\]。\end{enumerate}}

当 \(k\) 为 Artin 环，且 \(G\) 为拓扑平坦形式 \(k\)-群时，函子 \(\mathbf H(G)\) 显然由其在 \(k\) 处的值 \[
  \mathbf H(G)=\mathbf H(G)(k)
\] 确定。也称 \(\mathbf H(G)\) 为 \(G\) 的\emph{（协变）双代数}。\footnote{\label{sga3:VIIB:nde:82}编者在此引入了下文将用到的记号 \(\mathcal H^{\mathrm{cocom}}_{\mathrm{flat}/k}\)。}以 \(\mathcal H_k\) 表示 \(k\)-Hopf 代数范畴，以 \[
  \mathcal H^{\mathrm{cocom}}_{\mathrm{flat}/k}
\] 表示其中由在 \(k\) 上平坦且余交换的 \(k\)-Hopf 代数构成的全子范畴。于是得到以下推论。

\medskip\noindent\textbf{推论。}\phantomsection \label{sga3:VIIB:2.2.1.cor}{\itshape 设 \(k\) 为 Artin 环。\begin{enumerate}\item[\textup{(i)}]函子 \[
  G\longmapsto\mathbf H(G),
  \qquad
  \mathbf H\longmapsto
  \mathscr G(\mathbf H)=\operatorname{Spf}^*(\mathbf H)
\] 定义拓扑平坦形式 \(k\)-群范畴与 \(\mathcal H^{\mathrm{cocom}}_{\mathrm{flat}/k}\) 之间的等价。

\item[\textup{(ii)}]此等价“与基变换交换”：若 \(k\to K\) 为 Artin 环态射，则 \[
  \mathbf H\bigl(G\widehat\otimes_kK\bigr)
  =
  \mathbf H(G)\otimes_kK
\]，且 \[
  \mathscr G\bigl(\mathbf H\otimes_kK\bigr)
  =
  \mathscr G(\mathbf H)\widehat\otimes_kK.
\]。\end{enumerate}}

另一方面，以 \[
  \mathcal H^{\mathrm{com}}_{\mathrm{flat}/k}
\] 表示 \(\mathcal H_k\) 中由在 \(k\) 上平坦且交换的 \(k\)-Hopf 代数构成的全子范畴。回顾函子 \[
  T\longmapsto\mathcal O(T)
\] 是从仿射 \(k\)-群概形范畴到 \(\mathcal H^{\mathrm{com}}_{\mathrm{flat}/k}\) 的反等价。

\subsubsection*{2.2.2}
\phantomsection
\label{sga3:VIIB:2.2.2}

为简便起见，假设 \(k\) 为 Artin 环，并设 \(G\) 为拓扑平坦形式 \(k\)-群。\footnote{\label{sga3:VIIB:nde:83}编者扩充了下文原文，以引入记号 \(D'(G)\) 和 \(D(K)\)；参见 [Ca62]，\S~14。}则 \(G\) 交换，当且仅当其仿射代数 \(\mathcal A(G)\) 的余乘法余交换，而这又等价于双代数 \(\mathbf H(G)\) 的乘法交换。此时 \(\mathbf H(G)\) 为交换且余交换的 Hopf 代数，在 \(k\) 上平坦。因此，若 \[
  D'(G)=\operatorname{Spec}\mathbf H(G),
\]，则 \(D'(G)\) 为交换 \(k\)-群概形，在 \(k\) 上仿射且平坦。

反之，若 \(T\) 为这样的 \(k\)-群概形，则其仿射代数 \(\mathcal O(T)\) 是在 \(k\) 上平坦的余交换 \(k\)-余代数范畴中的交换群对象。因此 \SGARefLink{sga3:VIIB:1.3.5.C}{1.3.5.C} 给出拓扑平坦形式 \(k\)-群 \(D(T)\)，定义为 \[
  D(T)
  =
  \operatorname{Spf}^*\mathcal O(T)
  =
  \operatorname{Spf}(\mathcal A),
  \qquad
  \mathcal A=\mathcal O(T)^*.
\]。由 \SGARefLink{sga3:VIIB:1.3.5.D}{1.3.5.D}，有典范同构 \[
  G=\operatorname{Spf}^*\mathbf H(G),
  \qquad
  \mathbf H(D(T))=\mathcal O(T).
\]。从而得到典范同构 \[
\begin{aligned}
  D(D'(G))
  &=
  \operatorname{Spf}^*\mathcal O(D'(G))
  =
  \operatorname{Spf}^*\mathbf H(G)
  =
  G,
  \\
  D'(D(T))
  &=
  \operatorname{Spec}\mathbf H(D(T))
  =
  \operatorname{Spec}\mathcal O(T)
  =
  T.
\end{aligned}
\]。此外，若以 \(k\text{-}\operatorname{Gr.}\) 表示 \(k\)-群概形范畴，则 \SGARefLink{sga3:VIIB:2.2.1}{推论 2.2.1} 给出函子性同构 \[
  \operatorname{Hom}_{\operatorname{Grf}/k}\bigl(G,D(T)\bigr)
  \simeq
  \operatorname{Hom}_{\mathcal H_k}
    \bigl(\mathbf H(G),\mathcal O(T)\bigr)
  \simeq
  \operatorname{Hom}_{k\text{-}\operatorname{Gr.}}
    \bigl(T,D'(G)\bigr).
\]。这就证明了下面的命题。

\medskip\noindent\textbf{命题（Cartier 对偶性）。}\phantomsection \label{sga3:VIIB:2.2.2.prop}{\itshape 设 \(k\) 为 Artin 环。\begin{enumerate}\item[\textup{(i)}]函子 \[
  G\longmapsto D'(G)=\operatorname{Spec}\mathbf H(G),
  \qquad
  T\longmapsto D(T)=\operatorname{Spf}^*\mathcal O(T)
\] 诱导交换拓扑平坦形式 \(k\)-群范畴 \(\mathcal{FC}_k\) 与交换仿射平坦 \(k\)-群概形范畴 \(\mathcal{AC}_k\) 之间的反等价。换言之，\(G\) 与 \(D'(G)\)，以及 \(T\) 与 \(D(T)\)，分别通过下列等式联系：\footnote{\label{sga3:VIIB:nde:84}以 \(\mathcal{FC}_k^f\)、\(\mathcal{AC}_k^f\) 分别表示 \(\mathcal{FC}_k\)、\(\mathcal{AC}_k\) 中由满足 \(\mathbf H(G)\)、\(\mathcal O(T)\) 为有限 \(k\)-模 \textup{（因而有限局部自由）} 的对象 \(G\)、\(T\) 构成的全子范畴。则 \(\mathcal{FC}_k^f\) 和 \(\mathcal{AC}_k^f\) 均与在 \(k\) 上有限且平坦的交换、余交换 Hopf \(k\)-代数范畴 \(\mathcal C\) 具有相同的对象。对应 \(G\mapsto\mathbf H(G)\) 为协变，对应 \(T\mapsto\mathcal O(T)\) 为反变。因此重新得到 VII\textsubscript{A} \SGARefLink{sga3:VIIA:3.3.1}{3.3.1} 中已经见过的范畴 \(\mathcal{AC}_k\) 的“Cartier 对偶性”。}\[
  \mathbf H(G)=\mathcal O(D'(G)),
  \qquad
  \mathcal O(T)=\mathbf H(D(T)).
\]

\item[\textup{(ii)}]此反等价“与基变换交换”：若 \(k\to K\) 为 Artin 环态射，则 \[
  D'\bigl(G\widehat\otimes_kK\bigr)
  =
  D'(G)\otimes_kK
\]，且 \[
  D\bigl(T\otimes_kK\bigr)
  =
  D(T)\widehat\otimes_kK.
\]。

\item[\textup{(iii)}]特别地，若 \(k\) 为域，则得到交换形式 \(k\)-群范畴与交换仿射 \(k\)-群概形范畴之间的反等价，且此反等价与基域扩张交换。\end{enumerate}}


